\subsection{Escopo físico e dinâmico do manipulador robótico}
A fase de aproximação tratada neste documento é uma aproximação curta (\emph{close rendezvous}), em que a captura ocorre pela ferramenta (\emph{tool capture}) enquanto o \emph{chaser} permanece à deriva em relação ao \emph{target}. Até o instante final de captura não há contato entre as espaçonaves; assim, as forças de contato não foram modeladas.

As espaçonaves operam em regime de microgravidade. No subsistema do manipulador robótico, as forças gravitacionais são desprezadas. No subsistema orbital relativo, a translação/rotação relativas são descritas pelas Equações de Hill \cite{hill1878}. Portanto, este documento modela apenas a atitude da base e o movimento interno do manipulador.

Não há atmosfera modelada, logo, não há arrasto aerodinâmico. Perturbações externas como torque magnético e pressão de radiação solar também não são consideradas. As forças e torques de controle aplicados ao corpo-base são tratados como entradas ideais no centro de massa (\emph{Center of Mass, CoM}) do \emph{chaser}, e também não foi modelado o sistema de propulsão.

Todos os corpos (juntas, elos, espaçonaves, base) são rígidos, sem elasticidade, deformações, vibrações, folgas ou ruído. Todos os elementos são ideais, mas com limites físicos de rotação, velocidade e aceleração, tratados pela lei de controle.

\subsection{Atuadores e sensores}
Os atuadores das juntas aplicam torques (ou forças, no caso de junta prismática) denotados por $\tau_i$. Sensores de junta fornecem posição e velocidade. Para a base, torques/forças são entradas de controle ideais (sem detalhamento de atuadores).

\subsection{Restrições operacionais}
Para evitar colisões, o movimento do manipulador robótico é restringido a um espaço de trabalho (\emph{workspace}) que não intercepte a base do \emph{chaser} nem o \emph{target}. Além disso, foram estabelecidos limites de deslocamento, velocidade e aceleração por junta:
\[
q_{i,min} \le q_i \le q_{i,max},\qquad
|\dot q_i| \le \dot q_{i,max},\qquad
|\ddot q_i| \le \ddot q_{i,max}.
\]
As quatro primeiras juntas ($J_1$–$J_4$) são revolutas e a quinta junta ($J_5$) é do tipo \emph{prismática}, dessa forma, obtém-se a configuração $4R1P$, com curso limitado ao seu comprimento útil.

\subsection{Incertezas e variações}
Foram consideradas incertezas fixas por elo e para o conjunto (massas $m_i$, posições de CoM e momentos de inércia $\mathbf{I}_i$). As variáveis de estado são os ângulos/deslocamentos de junta, suas velocidades e acelerações.

\subsection{Critérios de validade do modelo}
O modelo é válido para taxas e acelerações típicas de operação do manipulador robótico. Não se cobrem fenômenos não rígidos como flexibilidade mecânicas e nem a flexibilidade do material, nem pequenas variações na taxa de atitude da base e também há ausência do aquecimento. A modelagem cinemática e dinâmica seguem \cite{craig2005}, enquanto que no contexto de microgravidade e \emph{rendezvous and docking/berthing} com base flutuante seguem \cite{fehse2003,wie2008space}.

\subsection{Limites operacinais estabelecidos}

A seguir são destacados os limites operacionais estabelecidos para o manipulador robótico: